% Part: first-order-logic
% Chapter: natural-deduction
% Section: proof-theoretic-notions

\documentclass[../../../include/open-logic-section]{subfiles}

\begin{document}

\iftag{FOL}
      {\olfileid[id]{fol}{ntd}{ptn}}
      {\olfileid[id]{pl}{ntd}{ptn}}

\olsection{Gagasan Teori-Bukti}

\begin{editorial}
Bagian ini menghimpun definisi relasi keterderivasian dan konsistensi
untuk deduksi alami.
\end{editorial}

\begin{explain}
Sebagaimana kita telah mendefinisikan sejumlah gagasan semantik penting
(validitas, perikutan, keterpenuhan), kini kita mendefinisikan
\emph{gagasan teori-bukti} yang bersesuaian. Gagasan-gagasan ini tidak
didefinisikan dengan mengacu pada terpenuhinya !!{sentence} di dalam
\iftag{FOL}{!!{structure}}{!!{valuation}}, melainkan dengan mengacu pada
!!{derivability} atau !!{nonderivability} !!{sentence} tertentu dari
yang lainnya. Penemuan penting menunjukkan bahwa gagasan-gagasan
teori-bukti tersebut berimpit dengan padanan semantiknya. Keberimpitan itu
adalah isi \emph{teorema kesahihan} dan \emph{teorema kelengkapan}.
\end{explain}


\begin{defn}[Teorema]
!!^a{sentence}~$!A$ adalah sebuah \emph{teorema} jika terdapat
!!a{derivation} untuk~$!A$ dalam deduksi alami yang semua asumsinya
!!{discharged}. Kita menulis $\Proves !A$ jika $!A$ merupakan teorema dan
$\Proves/ !A$ jika bukan.
\end{defn}

\begin{defn}[!!^{derivability}]
Kita mengatakan bahwa !!a{sentence} $!A$ \emph{!!{derivable} dari} suatu
himpunan !!{sentence}~$\Gamma$, dan menulis $\Gamma \Proves !A$, jika
terdapat !!a{derivation} dengan kesimpulan~$!A$ yang setiap asumsinya
entah !!{discharged} atau berada dalam~$\Gamma$. Jika $!A$ tidak
!!{derivable} dari $\Gamma$, kita menulis $\Gamma \Proves/ !A$.
\end{defn}

\begin{defn}[Konsistensi]
Suatu himpunan !!{sentence}~$\Gamma$ \emph{tak konsisten} jika dan hanya
jika $\Gamma \Proves \lfalse$. Jika $\Gamma$ tidak tak konsisten, yakni
jika $\Gamma \Proves/ \lfalse$, kita menyebutnya \emph{konsisten}.
\end{defn}

\begin{prop}[Refleksivitas]
\ollabel{prop:reflexivity}
Jika $!A \in \Gamma$, maka $\Gamma \Proves !A$.
\end{prop}

\begin{proof}
Asumsi $!A$ sendiri merupakan !!a{derivation} untuk~$!A$; semua asumsi
yang !!{undischarged} di dalamnya (yakni $!A$) berada dalam~$\Gamma$.
\end{proof}


\begin{prop}[Monotonisitas]
\ollabel{prop:monotonicity}
Jika $\Gamma \subseteq \Delta$ dan $\Gamma \Proves !A$, maka $\Delta
\Proves !A$.
\end{prop}

\begin{proof}
Setiap !!{derivation} untuk $!A$ dari $\Gamma$ juga merupakan
!!a{derivation} untuk $!A$ dari~$\Delta$.
\end{proof}

\begin{prop}[Transitivitas]
\ollabel{prop:transitivity}
Jika $\Gamma \Proves !A$ dan $\{!A\} \cup \Delta \Proves
!B$, maka $\Gamma \cup \Delta \Proves !B$.
\end{prop}

\begin{proof}
Jika $\Gamma \Proves !A$, terdapat !!a{derivation}~$\delta_0$ untuk~$!A$
yang semua asumsi yang !!{undischarged} di dalamnya berada dalam~$\Gamma$.
Jika
$\{!A\} \cup \Delta \Proves !B$, maka terdapat
!!a{derivation}~$\delta_1$ untuk~$!B$ yang semua asumsi yang
!!{undischarged} di dalamnya berada dalam~$\{!A\} \cup \Delta$. Sekarang
perhatikan:
\begin{prooftree}
  \AxiomC{$\Delta, \Discharge{!A}{1}$}
  \RightLabel{$\delta_1$}
  \DeduceC{$!B$}
  \DischargeRule{\Intro{\lif}}{1}
  \UnaryInfC{$!A \lif !B$}
  \AxiomC{$\Gamma$}
  \RightLabel{$\delta_0$}
  \DeduceC{$!A$}
  \RightLabel{\Elim{\lif}}
  \BinaryInfC{$!B$}
\end{prooftree}
Semua asumsi !!{undischarged} kini berada dalam $\Gamma \cup \Delta$,
sehingga hal ini menunjukkan $\Gamma \cup \Delta \Proves !B$.
\end{proof}

Jika $\Gamma = \{!A_1, !A_2, \ldots, !A_k\}$ merupakan himpunan
berhingga, kita dapat menggunakan notasi sederhana
$!A_1,!A_2,\ldots,!A_k \Proves !B$ untuk $\Gamma \Proves !B$; khususnya,
$!A \Proves !B$ berarti bahwa $\{!A\} \Proves !B$.

Perhatikan bahwa jika $\Gamma \Proves !A$ dan $!A \Proves !B$, maka
$\Gamma \Proves !B$. Diperoleh pula bahwa jika $!A_1, \dots, !A_n
\Proves !B$ dan $\Gamma \Proves !A_i$ untuk setiap~$i$, maka
$\Gamma \Proves !B$.

\begin{prop}
\ollabel{prop:incons}
Hal-hal berikut ekuivalen.
    \begin{enumerate}
      \item \( \Gamma \) tak konsisten.
      \item \( \Gamma \Proves {!A} \) untuk setiap !!{sentence}~\( {!A} \).
      \item \( \Gamma \Proves {!A} \) dan \( \Gamma \Proves \lnot {!A} \) untuk suatu !!{sentence}~\( {!A} \).
    \end{enumerate}
\end{prop}

\begin{proof}
Latihan.
\end{proof}

\tagprob{FOL}
\begin{prob}
Buktikan \olref[fol][ntd][ptn]{prop:incons}.
\end{prob}
\tagendprob

\tagprob{notFOL}
\begin{prob}
Buktikan \olref[pl][ntd][ptn]{prop:incons}.
\end{prob}
\tagendprob

\begin{prop}[Kekompakan]
  \ollabel{prop:proves-compact}
  \begin{enumerate}
  \item Jika $\Gamma \Proves !A$, maka terdapat himpunan bagian berhingga
    $\Gamma_0 \subseteq \Gamma$ sedemikian sehingga $\Gamma_0 \Proves !A$.
  \item Jika setiap himpunan bagian berhingga dari~$\Gamma$ konsisten,
    maka $\Gamma$ konsisten.
  \end{enumerate}
\end{prop}

\begin{proof}
  \begin{enumerate}
    \item Jika $\Gamma \Proves !A$, maka terdapat
      !!a{derivation}~$\delta$ untuk~$!A$ dari~$\Gamma$. Misalkan
      $\Gamma_0$ adalah himpunan asumsi !!{undischarged} dari~$\delta$.
      Karena setiap !!{derivation} berhingga, $\Gamma_0$ hanya dapat
      memuat berhingga banyak !!{sentence}. Jadi, $\delta$ merupakan
      !!a{derivation} untuk~$!A$ dari suatu $\Gamma_0 \subseteq \Gamma$
      yang berhingga.
    \item Ini adalah kontraposisi dari (1) untuk kasus khusus
      $!A \ident \lfalse$.
  \end{enumerate}
\end{proof}

\end{document}
